%- - - - - - - - - - - - - - - - - - - - - - - - - - - - - - - - - - - - - - - -
\subsubsection{Intraday Power Load Conventions}
\label{sss:intraday_power_load_conventions}
%- - - - - - - - - - - - - - - - - - - - - - - - - - - - - - - - - - - - - - - -

A node with name \lstinline!IntradayPowerLoad! is used to store reusable intraday power load profiles. These
conventions can be referenced from \lstinline!PowerLoadProfileReference! in an
\lstinline!IntradayPowerFloatingLegData! node. The same \lstinline!PowerLoadProfileData! structure can also be given
inline in the leg itself.

\begin{listing}[H]
\begin{minted}[fontsize=\footnotesize]{xml}
<IntradayPowerLoad>
  <Id>...</Id>
  <PowerLoadProfileData>
    ...
  </PowerLoadProfileData>
</IntradayPowerLoad>
\end{minted}
\caption{Intraday power load conventions}
\label{lst:intraday_power_load_conventions}
\end{listing}

The meanings of the elements in this node are as follows:
\begin{itemize}
\item \lstinline!Id!: Unique identifier of the load profile convention.
\item \lstinline!PowerLoadProfileData!: Load profile definition used to convert daily-average prices into load-weighted intraday prices.
\end{itemize}

\lstinline!PowerLoadProfileData! supports two mutually exclusive representations:
\begin{itemize}
\item \lstinline!ExplicitDates!: Explicit load profiles by effective date. Profiles are forward-flat filled, i.e. if a delivery date has no profile, the latest profile with effective date less than or equal to the delivery date is used.
\item \lstinline!BusinessDayRules!: Rule-based load profiles with a calendar and separate business-day and non-business-day profiles. Rules are also forward-flat filled by effective date.
\end{itemize}

In both cases, load is given by one or more \lstinline!LoadFactor! elements. Each \lstinline!LoadFactor! has required
attributes \lstinline!from! and \lstinline!to!, an optional \lstinline!unit! attribute and an optional
\lstinline!dst! flag. If \lstinline!unit! is omitted, seconds are assumed. The element value is the load in MW over
the interval. The \lstinline!dst! flag is intended for the extra hour between 02:00 and 03:00 on a daylight-saving
backward-shift date.

\begin{listing}[h!]
\begin{minted}[fontsize=\footnotesize]{xml}
<PowerLoadProfileData>
  <ExplicitDates>
    <LoadProfileDatum>
      <Date>2026-06-25</Date>
      <LoadFactors>
        <LoadFactor from="1" to="24" unit='HOUR'>1.0</LoadFactor>
      </LoadFactors>
    </LoadProfileDatum>
  </ExplicitDates>
</PowerLoadProfileData>

<PowerLoadProfileData>
  <BusinessDayRules>
    <LoadProfileBusinessDayRule>
      <Date>2026-01-01</Date>
      <Calendar>US-NERC</Calendar>
      <BusinessDayLoadFactors>
        <LoadFactor from="8" to="20" unit='HOUR'>1.0</LoadFactor>
      </BusinessDayLoadFactors>
      <NonBusinessDayLoadFactors>
        <LoadFactor from="1" to="24" unit='HOUR'>0.6</LoadFactor>
      </NonBusinessDayLoadFactors>
    </LoadProfileBusinessDayRule>
  </BusinessDayRules>
</PowerLoadProfileData>
\end{minted}
\caption{Power load profile data examples}
\label{lst:intraday_power_load_profile_data}
\end{listing}

In the first example, the profile is given explicitly for the date 2026-06-25. It defines a constant load of 1 MW
from hour 1 to hour 24 on that date. If no later explicit profile is provided, this profile is also used for later
 delivery dates until it is replaced by another explicit dated profile.

In the second example, the profile is rule-based from 2026-01-01 onward. The \lstinline!Calendar! determines whether a
 delivery date is treated as business day or non-business day. On business days in \lstinline!US-NERC!, the active load
profile is 1 MW from 08:00 to 20:00. On non-business days, the active load profile is 0.6 MW from 01:00 to 24:00.
As with explicit profiles, the rule is forward-flat filled and remains in force until a later rule date is provided.

Pricing uses the load profile to convert a daily-average price into a load-weighted daily price per MWh.

On a daylight-saving forward date, the missing hour between 02:00 and 03:00 is excluded from both the delivered MWh
and the load-weighted shape average. For example, assume a load profile of 50 MW from 01:00 to 08:00 and a
daily-average price of 100 USD/MWh. On a normal day, the delivered energy is $50 \times 7 = 350$ MWh. On a forward
day, the 02:00--03:00 hour does not exist, so the delivered energy becomes
$$
50 \times 6 = 300 \text{ MWh}.
$$